\section{Corrected polarized Einstein--scalar benchmark}
\label{metric}
Write $\tau=-\log t$, take $8\pi G=1$, and let $\varphi$ denote
the canonical massless scalar. Use $\phi=\sqrt{2}\varphi$ for the
equal-coefficient field variable. With positive metric quadratures
the metric is
\begin{align}
 ds^2={}&e^{(\tau-\lambda)/2}(-e^{-2\tau}d\tau^2+d\theta^2)\notag\\
   &+e^{-\tau}(e^Pdx^2+e^{-P}dy^2),\\
 \lambda_\tau={}& P_\tau^2+\phi_\tau^2
                 +e^{-2\tau}(P_\theta^2+\phi_\theta^2),\\
 \lambda_\theta={}& 2(P_\tau P_\theta+\phi_\tau\phi_\theta).
 \label{constraint}
\end{align}
The original preprint used the opposite sign for $\lambda$ in the
metric and a different scalar density normalization \cite{earlier}.
In the convention with $e^{(\tau+\lambda)/2}$, as in Lim,
both quadratures change sign \cite[Eqs.~(1), (4)--(5)]{lim}.
With normal pointing toward increasing $\tau$, the lapse and
geometric rates are
\begin{align}
 N&=e^{(-\lambda-3\tau)/4},\qquad
 H_\theta=\frac{1-\lambda_\tau}{4N},\\
 H_x&=\frac{P_\tau-1}{2N},\qquad H_y=\frac{-P_\tau-1}{2N},\\
 \Theta&=-\frac{3+\lambda_\tau}{4N},\\
 \rho&=\frac{\phi_\tau^2+e^{-2\tau}\phi_\theta^2}{4N^2}.
\end{align}
Let $H=\Theta/3$, $\Omega=\rho/(3H^2)$,
$\Sigma^2=\sum_i(H_i-H)^2/(6H^2)$, and $s_P=(H_x-H_y)/\Theta$.
The comparison uses these quantities on equal areal slices and the
dimensionless rate $\Theta e^{-\tau/2}$. A separate relative difference
in $\rho^{-1/2}$ compares density scales on the same slices.
If the unit normal is reversed, $\Theta$ changes sign and the
three normalized shape quantities stay fixed.

Among allowed asymptotic fields $P=v_P\tau+P_\infty+o(1)$,
$\phi=v_\phi\tau+\phi_\infty+o(1)$, assume differentiable asymptotics
of sufficient order to integrate the quadratures. The periodic
compatibility condition is
$\int_{\T}(v_PP_\infty'+v_\phi\phi_\infty')\,d\theta=0$.
With $q_0=v_P^2+v_\phi^2$, one has
$\lambda=q_0\tau+\lambda_\infty+o(1)$ and
$\lambda_\infty'=2(v_PP_\infty'+v_\phi\phi_\infty')$.
The integration constant $c_\lambda=\lambda_\infty(0)$ is
independent on the fixed marked torus \cite[Eq.~(3.19)]{beyer}. The shift
$\lambda\mapsto\lambda+c$ rescales $N$ and the transverse proper
length by $e^{-c/4}$, and normal rates by $e^{c/4}$.
These factors cancel from the normalized shape score.
The score also misses a constant shift of $P$, which changes
the two marked orbit lengths in reciprocal directions.
The full-state reconstruction of the marked metric requires
retaining its actual field data and integration constant.

The smooth production example has
$P=0.6\cos\theta$, $P_\tau=0.4+0.3\cos 2\theta$,
$\phi=0.5\sin\theta$, and $\phi_\tau=0.8-0.2\cos 2\theta$
at $\tau=0$, with $\lambda(0,0)=0$.
The reconstruction fits the two fields from the slices
$\tau_f-4$ and $\tau_f-2$ and estimates
$c_\lambda$ from $\lambda(\tau_f-2,0)-\widehat q_0(0)(\tau_f-2)$.
This supplies one pre-endpoint metric measurement in addition to
the field samples. Truth at the restart surface is withheld
from the main estimator. Centered second-order spatial differences
and RK4 evolve both branches.

At $N=1024$ the maximum geometric comparison distances on
$\tau=2,4$, for $\tau_f=8,12,16,24$, were
$4.759\times10^{-4}$, $1.060\times10^{-5}$,
$9.786\times10^{-6}$, and $9.786\times10^{-6}$ respectively.
At depths 16 and 24 each resolution doubling from 256 to 512 and
1024 lowered the error by about four. At depth 8 the refinement
sequence approached a finite-window error floor. Setting the estimated
constant to zero while retaining the four field functions produces
a score of about $0.232$ at the finest grid across all four depths.
In the analytic multimode Bessel benchmark the depth-16 fine-grid
scores were $1.3344\times10^{-5}$ both when the limiting coefficients
were supplied and when pre-endpoint holdout was used. The analytic
arm's finite-time metric truth uses the constrained quadrature of
the known fields. These results verify the declared outputs in the
chosen family and areal slicing.

\section{A finite-window nonlinear inverse}
\label{nonlinear}
For vacuum unpolarized Gowdy let $V=P_\tau$ and $R=Q_\tau$.
Given $P(f-\Delta),P(f),Q(f),R(f)$, the earlier method initialized
$V$ at $f$ to the trailing secant $v=(P(f)-P(f-\Delta))/\Delta$.
Its stored values were
\[
 (v, P(f)-fv, Q(f)+R(f)/(2v_{\rm safe}), e^{2P(f)}R(f)).
\]
At a known depth the original code's clipping convention makes
those values algebraically equivalent to the observations used
here. Thus the secant restart copies $P,Q,R$ at the endpoint.
It changes the missing $V$ component. The change in backward
field history measures the performance of the restart rule.

A finite-window inverse estimates $V$ by integrating the reduced
equations back to $f-\Delta$ from a trial endpoint velocity $V_n$.
If $P_{\mathrm{trial}}(f-\Delta;V_n)$ is the resulting field, the
iteration is
\[
 V_{n+1}=V_n+
 \frac{P_{\mathrm{trial}}(f-\Delta;V_n)-P(f-\Delta)}{\Delta}.
\]
Only the specified observations are used in this iteration.
The endpoint truth velocity and earlier comparison times are
reserved for evaluation. To see why the iteration can work, write
$P_{\mathrm{trial}}(f-\Delta;V)=P(f)-\Delta V+J(V)$, where
$J(V)=\int_0^\Delta(\Delta-u)P_{\tau\tau}(f-u;V)\,du$.
For a fixed finite-dimensional discretization, a state-flow
Lipschitz bound $e^{Lu}$ and acceleration Lipschitz bound $L_F$
make $J/\Delta$ Lipschitz with constant at most
$\frac12\Delta L_F e^{L\Delta}$ on an admissible ball. An
invariant ball and a bound below one would give local convergence.
We do not use this sufficient bound as a certificate for the
simulations.

At a homogeneous polarized vacuum background,
a finite Fourier Galerkin system gives a further analytic check.
When the initial $Q$ and $Q_t$ perturbations vanish, the
linearized equation for the unmeasured $P_t$ component
is the polarized equation. If all the retained two-time
multipliers are nonzero, the derivative of the map from that
velocity to the second field snapshot is invertible.
The finite-dimensional inverse function theorem then gives
a local inverse of this reduced observation map.
The neighborhood and constants can depend on the bandwidth.

The exact true-spike family from Lim provides a local
known-answer control \cite{lim,lim2}. The corrected form for
$Q$ already occurs in Appendix A of \cite{lim2}. For fixed
time, $e^{-P(x)}$ is quadratic about the spike center.
The center can be estimated from the vertex using sampled
$P$ alone. Evaluating the two fitted quadratics there gives
the finite-window center slope and fixes the family velocity
parameter. Center values of $P$ and $Q$ then determine the
remaining family parameters. This procedure uses information
available in the original tuple and the formulas of a known
family. An off-cell center can be recovered without providing
its grid shift. The calculation is local in $x$ and does not
identify a general periodic feature as a Lim solution.

The original production example had $P=R=0$,
$V=2+1.5\cos\theta$, and $Q=\cos\theta$ at $\tau=0$.
The ordinary secant and the inverse receive identical observations
at $f-2,f$ with $f=12$ or $18$. We retain the original mask at
depth 18, while also scoring the full periodic domain.
For two states, let $d_H$ be hyperbolic distance in the $(Q,e^{-P})$
upper half-plane, $s=(V^2+e^{2P}R^2)^{1/2}$, and
$E=s^2+e^{-2\tau}(P_\theta^2+e^{2P}Q_\theta^2)$.
The score is the sampled supremum of
$\bigl(d_H^2+(\Delta s)^2+[(\Delta E)/(1+|E_{\mathrm{ref}}|)]^2\bigr)^{1/2}$.
It involves the reduced fields and their derivatives.
The comparison omits the value of the metric potential.

At $f=18$ the following scores are measured at the held-out
slice $\tau=2$. The legacy forward stencil is second order.
The two inverse columns use a fourth-order restart. The
privileged case supplies the endpoint velocity from the reference.
\begin{center}
\begin{tabular}{rccccc}
\toprule
 & \multicolumn{2}{c}{Secant, legacy stencil} &
 \multicolumn{2}{c}{Window inverse, fourth order} & Privileged \\
$N$ & Mask & Full circle & Mask & Full circle & Full circle\\
\midrule
1024 & $0.82012$ & $2.73288$ & $0.04666$ & $0.04666$ & $0.04734$\\
2048 & $0.81689$ & $3.09459$ & $0.01186$ & $0.01186$ & $0.01203$\\
4096 & $0.82215$ & $3.03402$ & $0.00298$ & $0.00298$ & $0.00302$\\
\bottomrule
\end{tabular}
\end{center}
At depth 18 the window fit residual divided by the observation
gap is below $10^{-10}$ at every tested grid. On the finest
grid the inverse gives $1.15\times10^{-3}$ and $1.55\times10^{-4}$
at $\tau=4$ and $8$. Depth 12 has no flagged component.
There the fine-grid full-circle score at $\tau=2$ falls from
$0.43207$ to $9.95\times10^{-5}$.
A same-stencil inverse is also recorded and tests inversion of its
own discretization. An analytic periodic geodesic wave-map control
provides snapshots independently of the solver. Its fourth-order
window inverse reduces the $\tau=2$, $N=1024$ score from
$0.02577$ to $1.62\times10^{-9}$.
For the two noiseless local Lim controls, the estimated half-cell
center errors were at most $1.45\times10^{-15}$ in coordinates.
Those controls test recovery within a specified family.

For a separate nonlinear check, evolve from $P=0$,
$V=2+1.4\cos\theta$, $Q=0.9\cos\theta+0.1\sin2\theta$,
and $R=0.025\cos3\theta$, with a periodic primitive of the
momentum source. Observe at $\tau=10,12$ using a fourth-order
spatial stencil. Reconstruct using the historical second-order
stencil and compare at $\tau=2$.
Apply the same reduced full-circle score.

\begin{center}
\begin{tabular}{rccc}
\toprule
$N$ & Secant & Window inverse & Exact endpoint velocity\\
\midrule
512  & $0.23690$ & $0.009585$ & $0.009876$\\
1024 & $0.23614$ & $0.002405$ & $0.002479$\\
\bottomrule
\end{tabular}
\end{center}
The matched-window score is close to that of a restart supplied
with the exact numerical endpoint velocity. On the finer grid,
adding an absolute error of $10^{-7}$ to the earlier $P$
observation in a single cosine mode gave gains of 5.17 and 8.36
in the comparison score at modes 1 and 128, respectively.
These are two directional measurements for this IVP. They
provide no infinite-band bound. The thin-feature classification
and the asymptotic regime of nonlinear unpolarized Gowdy
remain questions distinct from finite-window reconstruction.


\section{Conclusions}
In the polarized problem, the same pair of field snapshots has
different regularity costs for different targets. A rational gap can
make the field reconstruction stable at specified slices while the
missing velocity remains sensitive to high-frequency errors.
Assigning a sampling time to each band removes that additional
loss in the energy-normalized velocity.
The nonlinear comparison shows that the large error in the
original finite-window test can be removed with the same inputs
by fitting the PDE evolution across that window.
A uniform nonlinear inverse estimate near true spikes remains
open. The analytic statements concern a fixed periodic marking and
the specified finite-time observation maps.

\appendix
\section{Derivation of the scalar and metric signs}
In areal time write the base metric as $h=e^{2\sigma}(-dt^2+d\theta^2)$,
where $\sigma=-(\log t+\lambda)/4$.
The two one-dimensional warped directions have logarithmic
factors $A=(\log t+P)/2$ and $B=(\log t-P)/2$.
The base components of the four-dimensional Ricci tensor are
$R_{ab}(h)-\nabla^h_a\nabla^h_b(A+B)-A_{,a}A_{,b}-B_{,a}B_{,b}$.
Its mixed component and the sum of the two diagonal components are
\begin{align*}
 R_{t\theta}&=-\lambda_\theta/(4t)-P_tP_\theta/2,\\
 R_{tt}+R_{\theta\theta}
   &=-\lambda_t/(2t)-(P_t^2+P_\theta^2)/2 .
\end{align*}
Apply $R_{ab}=\partial_a\varphi\,\partial_b\varphi$ to obtain
\eqref{constraint}. The remaining base difference follows by
differentiating those quadratures and using the wave equations.
The two orbit equations reduce to \eqref{wave} for $P$.
The scalar d'Alembertian gives the same equation for $\varphi$.
As a check on homogeneous slices, their intrinsic scalar curvature
vanishes, so $G(n,n)=H_\theta H_x+H_xH_y+H_yH_\theta$.
With $P_\tau=p$, $\phi_\tau=v$ and $\lambda_\tau=p^2+v^2$
the expression reduces to $v^2/(4N^2)$. The vacuum example $p=1,v=0$
and the stiff isotropic example $p=0,v=\sqrt{3}$ pass this constraint.
